\documentclass[11pt]{article}
\usepackage[margin=1in]{geometry}
\usepackage{amsmath, amssymb, amsthm}
\usepackage{mathtools}
\usepackage{hyperref}

\newtheorem{theorem}{Theorem}
\newtheorem{definition}{Definition}
\newtheorem{proposition}{Proposition}

\title{Expectation}
\author{math-foundations / 25-expectation}
\date{}

\begin{document}
\maketitle

\section{Definitions}

\begin{definition}[Expectation, discrete]
Let $X$ be a discrete random variable on a probability space
$(\Omega, \mathcal{F}, P)$ taking values $\{k_i\}_{i \ge 1} \subset \mathbb{R}$
with probability mass function $p_X(k) = P(X = k)$. If
$\sum_{i} |k_i|\, p_X(k_i) < \infty$, the \emph{expectation} of $X$ is
\[
E[X] \;=\; \sum_{i} k_i\, p_X(k_i).
\]
\end{definition}

\begin{definition}[Expectation, continuous]
Let $X$ be a continuous random variable with probability density function
$f_X$. If $\int_{-\infty}^{\infty} |x|\, f_X(x)\, dx < \infty$, then
\[
E[X] \;=\; \int_{-\infty}^{\infty} x\, f_X(x)\, dx.
\]
\end{definition}

\begin{definition}[Expectation, Lebesgue]
For an integrable measurable $X : \Omega \to \mathbb{R}$,
\[
E[X] \;=\; \int_{\Omega} X(\omega)\, dP(\omega).
\]
This is constructed first for simple functions, extended to non-negative
measurable functions via monotone approximation, and to signed functions
via $X = X^+ - X^-$.
\end{definition}

\section{Linearity of Expectation}

\begin{theorem}[Linearity of Expectation]
\label{thm:linearity}
Let $X$ and $Y$ be random variables on a common probability space with
$E[|X|] < \infty$ and $E[|Y|] < \infty$, and let $a, b \in \mathbb{R}$. Then
\[
E[aX + bY] \;=\; a\,E[X] + b\,E[Y].
\]
Linearity holds whether or not $X$ and $Y$ are independent.
\end{theorem}

\begin{proof}[Proof (discrete case)]
Assume $X$ and $Y$ are discrete random variables with joint pmf
$p_{X,Y}(x, y) = P(X = x, Y = y)$, and that
$\sum_{x,y} |x|\, p_{X,Y}(x,y) < \infty$ and likewise for $Y$. Let
$Z = aX + bY$.

By LOTUS applied to the joint pmf,
\begin{align*}
E[Z] \;=\; E[aX + bY]
     &= \sum_{x} \sum_{y} (a x + b y)\, p_{X,Y}(x, y) \\
     &= a \sum_{x} \sum_{y} x\, p_{X,Y}(x, y)
        \;+\; b \sum_{x} \sum_{y} y\, p_{X,Y}(x, y).
\end{align*}
Absolute convergence (from the integrability assumptions) permits us to
rearrange the double sum and apply Fubini for series. Marginalising the
joint pmf gives $\sum_{y} p_{X,Y}(x, y) = p_X(x)$ and
$\sum_{x} p_{X,Y}(x, y) = p_Y(y)$, so
\begin{align*}
E[Z] &= a \sum_{x} x \left( \sum_{y} p_{X,Y}(x, y) \right)
       + b \sum_{y} y \left( \sum_{x} p_{X,Y}(x, y) \right) \\
     &= a \sum_{x} x\, p_X(x) + b \sum_{y} y\, p_Y(y) \\
     &= a\, E[X] + b\, E[Y].
\end{align*}
No use of independence was required: the joint pmf $p_{X,Y}$ carries any
dependence structure, and it factors out of each single-variable sum only
through its marginals.
\end{proof}

\begin{proposition}[Jensen's inequality]
Let $\phi : \mathbb{R} \to \mathbb{R}$ be convex and let $X$ be an integrable
random variable with $E[|\phi(X)|] < \infty$. Then
\[
\phi\!\big(E[X]\big) \;\le\; E\!\big[\phi(X)\big].
\]
\end{proposition}

\begin{proposition}[LOTUS]
For measurable $g : \mathbb{R} \to \mathbb{R}$ such that $g(X)$ is integrable,
\[
E[g(X)] \;=\; \int_{\mathbb{R}} g(x)\, f_X(x)\, dx
\quad \text{(continuous)},
\qquad
E[g(X)] \;=\; \sum_{k} g(k)\, p_X(k)
\quad \text{(discrete)}.
\]
\end{proposition}

\section{Examples}

\paragraph{Bernoulli.} If $X \sim \mathrm{Bernoulli}(p)$, then
$E[X] = 1 \cdot p + 0 \cdot (1-p) = p$.

\paragraph{Uniform $[a,b]$.} If $X \sim \mathrm{Uniform}[a,b]$, then
\[
E[X] = \int_a^b \frac{x}{b-a}\, dx = \frac{a+b}{2}.
\]

\paragraph{Binomial.} If $X \sim \mathrm{Binomial}(n, p)$, write
$X = \sum_{i=1}^{n} X_i$ with $X_i \sim \mathrm{Bernoulli}(p)$ iid. By
Theorem~\ref{thm:linearity}, $E[X] = \sum_{i=1}^{n} E[X_i] = np$.

\end{document}
